\section{Case Identifier Grammar}

Table~\ref{tab:id_examples} decodes representative identifiers so that fleet size, geometry, Central Node fraction, and phasing can be traced from a result row back to its architecture.

\begin{lstlisting}
T{satellites:03d}_P{planes:02d}_H{altitude_km:04d}_
I{inclination_x10:04d}_CN{fraction:03d}_F1
\end{lstlisting}

\begin{table}[htbp]
\centering
\caption{Example case identifiers.}
\label{tab:id_examples}
\begin{tabularx}{\textwidth}{L{0.50\textwidth}Y}
\toprule
Identifier & Interpretation \\
\midrule
\texttt{T060\_P03\_H0500\_I0600\_CN000\_F1} & 60 satellites, 3 planes, 500 km, 60 deg, no central nodes \\
\texttt{T120\_P10\_H0800\_I0986\_CN030\_F1} & 120 satellites, 10 planes, 800 km, 98.6 deg, 30\% CN \\
\texttt{T180\_P10\_H1000\_I0600\_CN100\_F1} & 180 satellites, 10 planes, 1000 km, 60 deg, all nodes central \\
\bottomrule
\end{tabularx}
\end{table}

\section{Final Campaign Frozen Controls}

Table~\ref{tab:final_controls_appendix} records the fixed physical and operational settings that are not part of the architecture sweep.

\begin{table}[htbp]
\centering
\caption{Core physical and operational controls of the final verified campaign.}
\label{tab:final_controls_appendix}
\begin{tabularx}{\textwidth}{L{0.38\textwidth}L{0.25\textwidth}Y}
\toprule
Control & Value & Note \\
\midrule
Campaign seed & 20260811 & Scenario seeds derived by SHA-256 \\
Orbital time step & 60 s & Discrete opportunity evaluation \\
Observation radius & 700 km & Fixed; no radius sweep in the final campaign \\
Ground station & Munich & 10-deg minimum elevation \\
Sat--sat rate & V19 UHF calculation & Effective cap near 19.2 kbit/s \\
Sat--ground rate & 100 kbit/s & Fixed window rate \\
Central link-budget boost & 1.5 & Physical-rate multiplier in link calculation \\
Nominal topology & \texttt{central\_only} & CN0 has ground uplink but no ISL \\
Nominal routing & targeted useful/deadline & Greedy temporal replay \\
Nominal hop/copy limits & 3 / 3 observer / 5 CN & Removed in unlimited diagnostic \\
Window utilization & 1.0 & Capacity multiplier before failure \\
\bottomrule
\end{tabularx}
\end{table}

The legacy field \path{sat_sat_data_rate_bps=250000} is present in configuration provenance but is not used by the active V19 link-budget calculation. It must not be reported as the satellite-to-satellite rate of the final verified campaign.

\section{Scenario Counts}

Table~\ref{tab:scenario_counts_appendix} reconciles the deterministic and seeded scenario families with the 539 rows declared for every architecture.

\begin{table}[htbp]
\centering
\caption{Scenario rows per architecture.}
\label{tab:scenario_counts_appendix}
\begin{tabular}{lrrr}
\toprule
Family & Levels & Replicates/level & Rows \\
\midrule
Core/policy/unlimited & 7 & 1 & 7 \\
Task size & 4 & 1 & 4 \\
Random satellite failure & 4 & 30 & 120 \\
Random central-node failure & 4 & 30 & 120 \\
Link-window outage & 4 & 30 & 120 \\
Ground-station outage & 4 & 30 & 120 \\
Targeted CN failure & 4 & 1 & 4 \\
Capacity degradation & 4 & 1 & 4 \\
Combined stress & 4 & 10 & 40 \\
\midrule
Total & & & 539 \\
\bottomrule
\end{tabular}
\end{table}

\section{Seed Derivation}

Table~\ref{tab:seed_scope} states which factors share a random draw and which receive a new draw, preserving paired comparisons across architectures and regions.

For every stochastic row, the seed material is the UTF-8 concatenation

\begin{lstlisting}
campaign_seed|case_id|scenario_id|replicate
\end{lstlisting}

The SHA-256 digest is calculated and its first eight bytes define the numeric random seed. This makes each replicate stable even when worker count, job order, or resume timing changes.

\begin{table}[htbp]
\centering
\small
\caption{Random-seed sharing scope in the completed campaign.}
\label{tab:seed_scope}
\begin{tabularx}{\textwidth}{L{0.32\textwidth}L{0.22\textwidth}Y}
\toprule
Comparison & Shared random draw? & Reason \\
\midrule
Same case/scenario/replicate after resume & Yes & Complete seed key is identical \\
Different replicate of one case/scenario & No & Replicate number changes \\
Different CN fraction or architecture & No & Case identifier changes \\
Different scenario identifier & No & Scenario identifier changes \\
California versus India & Only if every seed-key field is identical & No cross-window sharing is assumed without an explicit paired contract \\
\bottomrule
\end{tabularx}
\end{table}

The current campaign therefore guarantees reproducibility, not common random numbers across architectures. Common random numbers can reduce variance only when the corresponding draws are deliberately shared and recorded across alternatives \citep{SloanUnwin1990}. A future paired variance-reduction run must therefore introduce a separate shared-draw identifier.

\section{Checkpoint and Completion Rules}

The runner writes scenario metrics in Parquet batches of 20 rows. It writes to a temporary path first and renames that path only after the batch is complete. On restart, the runner reads the scenario keys already stored and calculates the missing keys. An interruption close to a batch boundary can still leave duplicate keys, so the post-processing applies one deterministic key-level de-duplication rule before validation.

A case is classified as complete only if \path{_SUCCESS.json} exists and every expected scenario identifier is either completed or explicitly marked not applicable. Readable rows without that marker are not promoted into the complete analysis population.

\section{Representative Parquet Reading Commands}

Python with pandas and PyArrow can inspect a partition without conversion:

The literal directory name \path{full891_fresh_complete} in the following command is a legacy repository identifier for the clean, validated campaign restart. It does not denote a separate simulation method, scenario family, or evidence population.

\begin{lstlisting}
from pathlib import Path
import pandas as pd

case_dir = Path("outputs/full891_fresh_complete/runs/California") / CASE_ID
metrics = pd.concat(
    [pd.read_parquet(p) for p in sorted(case_dir.glob("case_metrics_part_*.parquet"))],
    ignore_index=True,
)
print(metrics.shape)
print(metrics[["scenario_id", "status",
               "observation_success_percent"]].head())
\end{lstlisting}

DuckDB provides an interactive alternative:

\begin{lstlisting}
SELECT scenario_id, status, observation_success_percent
FROM read_parquet('case_metrics_part_*.parquet')
ORDER BY scenario_id
LIMIT 20;
\end{lstlisting}

The compact CSV tables in this report package are intended for quick inspection and plotting. They are derived evidence and do not replace the raw Parquet archive.
